% tensor_cube.tex
% =====================================================================
% The 3D bucket-membership tensor with one window slice $M^{(j)}$
% highlighted, one specific bucket-column called out ($k = 3$, where
% the sparse pattern lands 5 nonzeros), and the pair-tree reduction
% of that bucket drawn IN PLACE — growing upward out of the slice
% along the $j$ axis. Single self-contained figure.
% Compile with `./build.sh` (or pdflatex tensor_cube.tex).
% =====================================================================
\documentclass[tikz,border=0pt]{standalone}
\usepackage{tikz}
\usepackage{tikz-3dplot}
\usepackage{amsmath}
\usetikzlibrary{calc,arrows.meta,positioning,backgrounds}

% Dark-mode palette — deep navy background, light strokes, saturated
% reds/golds for the highlighted bucket + tree.
\definecolor{bgDark}      {RGB}{15,25,45}
\definecolor{axisColor}   {RGB}{210,220,240}
\definecolor{cubeEdge}    {RGB}{120,160,220}
\definecolor{cubeFaint}   {RGB}{90,110,140}
\definecolor{slicePlane}  {RGB}{130,175,230}
\definecolor{oneCell}     {RGB}{220,110,110}
\definecolor{oneCellEdge} {RGB}{255,160,160}
\definecolor{fgText}      {RGB}{170,180,195}
\definecolor{colhi}       {RGB}{70,25,35}
\definecolor{colstroke}   {RGB}{255,130,130}
\definecolor{nodeFill}    {RGB}{60,50,25}
\definecolor{nodeEdge}    {RGB}{240,200,100}
\definecolor{carryFill}   {RGB}{55,60,75}

\tdplotsetmaincoords{68}{120}

\begin{document}
\begin{tikzpicture}[font=\normalsize, >={Latex[length=2.6mm, width=2.4mm]}]

\begin{scope}[
    tdplot_main_coords,
    scale=1.0,
    every node/.style={font=\small},
    axis/.style       ={->, thick, axisColor, line cap=round},
    edgeFront/.style  ={cubeEdge, semithick},
    edgeBack/.style   ={cubeFaint, thin, dashed},
    sliceEdge/.style  ={slicePlane, thick},
    nonzeroCell/.style={fill=oneCell, draw=oneCellEdge, very thin, fill opacity=0.9},
    hiCell/.style     ={fill=colstroke, draw=colstroke!70!black, very thin, fill opacity=0.95},
]

    \def\Ni{20}
    \def\Nk{10}
    \def\Nj{4}
    \def\jSlice{1}
    \def\khl{3}                  % bucket-column highlighted on the slice

    \pgfmathtruncatemacro{\NiMinusOne}{\Ni-1}
    \pgfmathtruncatemacro{\NkMinusOne}{\Nk-1}
    \pgfmathtruncatemacro{\NjMinusOne}{\Nj-1}

    % 5 of these (at i = 0, 3, 8, 11, 14) land in k = 3; those are the
    % 5 leaves of the pair-tree drawn above the slice.
    \def\sparsePattern{%
        0/3, 1/8, 2/1, 3/3, 4/0, 5/7, 6/2, 7/9, 8/3, 9/6,%
        10/1, 11/3, 12/7, 13/0, 14/3, 15/2, 16/8, 17/6, 18/9, 19/4}

    % --- Cube wireframe (back face).
    \draw[edgeBack]
        (0,0,\Nj) -- (\Ni,0,\Nj) -- (\Ni,\Nk,\Nj) -- (0,\Nk,\Nj) -- cycle;
    \foreach \xc/\yc in {0/0, \Ni/0, \Ni/\Nk, 0/\Nk}
        \draw[edgeBack] (\xc,\yc,0) -- (\xc,\yc,\Nj);

    % --- (j, k) grid on the i = 0 face.
    \foreach \kk in {1,...,\NkMinusOne}
        \draw[fgText, thin, densely dotted, opacity=0.55]
            (0, \kk, 0) -- (0, \kk, \Nj);
    \foreach \jj in {1,...,\NjMinusOne}
        \draw[fgText, thin, densely dotted, opacity=0.55]
            (0, 0, \jj) -- (0, \Nk, \jj);

    % --- Highlighted slice $M^{(j)}$.
    \fill[slicePlane, opacity=0.10]
        (0,0,\jSlice) -- (\Ni,0,\jSlice) -- (\Ni,\Nk,\jSlice) -- (0,\Nk,\jSlice) -- cycle;
    \draw[sliceEdge]
        (0,0,\jSlice) -- (\Ni,0,\jSlice) -- (\Ni,\Nk,\jSlice) -- (0,\Nk,\jSlice) -- cycle;

    \foreach \xc in {1,...,\NiMinusOne}
        \draw[slicePlane, very thin, opacity=0.20] (\xc,0,\jSlice) -- (\xc,\Nk,\jSlice);
    \foreach \yc in {1,...,\NkMinusOne}
        \draw[slicePlane, very thin, opacity=0.20] (0,\yc,\jSlice) -- (\Ni,\yc,\jSlice);

    % --- Highlighted bucket column k = khl on the slice.
    \fill[colhi, opacity=0.55]
        (0, \khl, \jSlice) -- (\Ni, \khl, \jSlice) --
        (\Ni, \khl+1, \jSlice) -- (0, \khl+1, \jSlice) -- cycle;
    \draw[colstroke, thick]
        (0, \khl, \jSlice) -- (\Ni, \khl, \jSlice) --
        (\Ni, \khl+1, \jSlice) -- (0, \khl+1, \jSlice) -- cycle;

    % --- Off-column nonzero cells (default red, behind tree).
    \foreach \ihit/\khit in \sparsePattern {
        \pgfmathtruncatemacro{\isHi}{ifthenelse(\khit == \khl, 1, 0)}
        \ifnum\isHi=0
            \fill[nonzeroCell]
                (\ihit,   \khit,   \jSlice) --
                (\ihit+1, \khit,   \jSlice) --
                (\ihit+1, \khit+1, \jSlice) --
                (\ihit,   \khit+1, \jSlice) -- cycle;
        \fi
    }
    % --- On-column nonzero cells (deeper red — they feed the tree).
    \foreach \ihit/\khit in \sparsePattern {
        \pgfmathtruncatemacro{\isHi}{ifthenelse(\khit == \khl, 1, 0)}
        \ifnum\isHi=1
            \fill[hiCell]
                (\ihit,   \khit,   \jSlice) --
                (\ihit+1, \khit,   \jSlice) --
                (\ihit+1, \khit+1, \jSlice) --
                (\ihit,   \khit+1, \jSlice) -- cycle;
        \fi
    }

    % --- Cube front face (drawn on top of slice fill).
    \draw[edgeFront]
        (0,0,0) -- (\Ni,0,0) -- (\Ni,\Nk,0) -- (0,\Nk,0) -- cycle;

    % =================================================================
    % Pair-tree growing upward (along +j) out of the highlighted column.
    % Leaves at j = zL0, climbing through three levels to a root at zL3.
    % All nodes live at k = khl + 0.5 (centred in the column); the i
    % coordinate of each interior node is the mid-point of the children.
    % Carry (P_14) keeps its i position across levels.
    % =================================================================
    % Tree levels live well above the j-axis tip (j = Nj + 1 = 5), so the
    % illustration sits above the cube and its label without intermixing.
    \def\zLeaf{9.0}      % leaves
    \def\zPair{10.2}
    \def\zQuad{11.4}
    \def\zRoot{12.6}     % root B_{j,k}

    % Leaf i-positions = i centres of the 5 nonzeros in column khl.
    \def\xA{0.5}  \def\xB{3.5}  \def\xC{8.5}  \def\xD{11.5} \def\xE{14.5}
    \pgfmathsetmacro{\xAB}{(\xA+\xB)/2}   % level-1 pair (A,B) mid
    \pgfmathsetmacro{\xCD}{(\xC+\xD)/2}   % level-1 pair (C,D) mid
    \pgfmathsetmacro{\xABCD}{(\xAB+\xCD)/2} % level-2 pair (AB,CD) mid
    \pgfmathsetmacro{\xRoot}{(\xABCD+\xE)/2} % level-3 root mid

    \tikzset{
        treeLeaf/.style ={circle, fill=oneCell, draw=oneCellEdge,
                          inner sep=0pt, minimum size=4.2pt},
        treeInode/.style={circle, fill=nodeFill, draw=nodeEdge, thick,
                          inner sep=0pt, minimum size=4.6pt},
        treeCarry/.style={circle, fill=carryFill, draw=nodeEdge, thick,
                          dashed, inner sep=0pt, minimum size=4.2pt},
        treeRoot/.style ={rectangle, rounded corners=2pt,
                          fill=nodeFill, draw=nodeEdge, thick,
                          text=axisColor, inner sep=3pt, font=\normalsize},
        treeEdge/.style ={nodeEdge, semithick, -},
        treeCarryEdge/.style={nodeEdge, semithick, dashed, -},
        feedEdge/.style ={colstroke, semithick, dotted},
    }

    % Cell-to-leaf feeders (short dotted lines lifting each cell up).
    \foreach \xc/\lbl in {\xA/$P_0$, \xB/$P_3$, \xC/$P_8$, \xD/$P_{11}$, \xE/$P_{14}$}
        \draw[feedEdge] (\xc, \khl+0.5, \jSlice) -- (\xc, \khl+0.5, \zLeaf);

    % Level 0 — leaves with labels (all placed above for consistency).
    \node[treeLeaf,label={[font=\footnotesize, text=axisColor]above:$P_0$}]
        (Ta) at (\xA, \khl+0.5, \zLeaf) {};
    \node[treeLeaf,label={[font=\footnotesize, text=axisColor]above:$P_3$}]
        (Tb) at (\xB, \khl+0.5, \zLeaf) {};
    \node[treeLeaf,label={[font=\footnotesize, text=axisColor]above:$P_8$}]
        (Tc) at (\xC, \khl+0.5, \zLeaf) {};
    \node[treeLeaf,label={[font=\footnotesize, text=axisColor]above:$P_{11}$}]
        (Td) at (\xD, \khl+0.5, \zLeaf) {};
    \node[treeLeaf,label={[font=\footnotesize, text=axisColor]above:$P_{14}$}]
        (Te) at (\xE, \khl+0.5, \zLeaf) {};

    % Level 1.
    \node[treeInode] (Tab)  at (\xAB, \khl+0.5, \zPair) {};
    \node[treeInode] (Tcd)  at (\xCD, \khl+0.5, \zPair) {};
    \node[treeCarry] (Te1)  at (\xE,  \khl+0.5, \zPair) {};
    \draw[treeEdge]      (Ta) -- (Tab);  \draw[treeEdge] (Tb) -- (Tab);
    \draw[treeEdge]      (Tc) -- (Tcd);  \draw[treeEdge] (Td) -- (Tcd);
    \draw[treeCarryEdge] (Te) -- (Te1);

    % Level 2.
    \node[treeInode] (Tabcd) at (\xABCD, \khl+0.5, \zQuad) {};
    \node[treeCarry] (Te2)   at (\xE,    \khl+0.5, \zQuad) {};
    \draw[treeEdge]      (Tab) -- (Tabcd);
    \draw[treeEdge]      (Tcd) -- (Tabcd);
    \draw[treeCarryEdge] (Te1) -- (Te2);

    % Level 3 — root $B_{j,k}$.
    \node[treeRoot] (Troot) at (\xRoot, \khl+0.5, \zRoot) {$B_{j,k}$};
    \draw[treeEdge] (Tabcd) -- (Troot);
    \draw[treeEdge] (Te2)   -- (Troot);

    % --- Axes (drawn last so they sit above the slice fill).
    \draw[axis] (0,0,0) -- (\Ni+1.0, 0, 0) coordinate (iTip);
    \draw[axis] (0,0,0) -- (0, \Nk+1.0, 0) coordinate (kTip);
    \draw[axis] (0,0,0) -- (0, 0, \Nj+1.0) coordinate (jTip);

    \node[font=\normalsize, text=axisColor, anchor=east, xshift=-4pt, inner sep=1.5pt]
        at (iTip) {$i$ \footnotesize ($n$ inputs, $2^{16}$–$2^{20}$)};
    \node[font=\normalsize, text=axisColor, anchor=west, xshift=4pt, inner sep=1.5pt]
        at (jTip) {$j$ \footnotesize ($T$ windows, $\approx 17$)};
    \node[font=\normalsize, text=axisColor, anchor=west, xshift=4pt, inner sep=1.5pt]
        at (kTip) {$k$ \footnotesize ($B_W$ buckets, $\sim 2^{14}$)};

    \node[font=\footnotesize, text=axisColor, anchor=north east]
        at (-0.05, -0.05, 0) {$0$};

    % --- Slice label.
    \node[slicePlane, font=\normalsize, anchor=west]
        (sliceLabel) at (-2.8, -1.2, \jSlice) {$M^{(j)}$};
    \draw[slicePlane, semithick, opacity=0.55]
        (sliceLabel.west) -- (0, 0, \jSlice);

    % --- Bucket-column label.
    \node[colstroke, font=\normalsize, anchor=north west]
        at (\Ni+0.3, \khl+0.5, \jSlice) {$k = \khl$};

\end{scope}

% Dark background — drawn on the background layer so it sits behind
% every other element while still covering the final bounding box.
\begin{scope}[on background layer]
  \fill[bgDark]
    ([shift={(-10pt,-10pt)}]current bounding box.south west)
    rectangle
    ([shift={( 10pt, 10pt)}]current bounding box.north east);
\end{scope}

\end{tikzpicture}
\end{document}
